% ========================================================
% BÀI TẬP TOÁN 5 - KẾT NỐI TRI THỨC VỚI CUỘC SỐNG
% Bài 33: Ôn tập diện tích, chu vi một số hình phẳng
% Định dạng: BTTeX Part III - Trắc nghiệm điền số ngắn
% ========================================================

% Câu 1
\begin{ex}
  Một hình vuông có cạnh $8\text{ cm}$. Chu vi hình vuông là bao nhiêu xăng-ti-mét?
  \shortans{32}
  \loigiai{
    $8 \times 4 = 32\text{ cm}$.
  }
\end{ex}

% Câu 2
\begin{ex}
  Một hình vuông có cạnh $8\text{ cm}$. Diện tích hình vuông là bao nhiêu xăng-ti-mét vuông?
  \shortans{64}
  \loigiai{
    $8 \times 8 = 64\text{ cm}^2$.
  }
\end{ex}

% Câu 3
\begin{ex}
  Một hình chữ nhật có chiều dài $15\text{ m}$, chiều rộng $8\text{ m}$. Chu vi hình chữ nhật là bao nhiêu mét?
  \shortans{46}
  \loigiai{
    $(15 + 8) \times 2 = 46\text{ m}$.
  }
\end{ex}

% Câu 4
\begin{ex}
  Một hình chữ nhật có chiều dài $15\text{ m}$, chiều rộng $8\text{ m}$. Diện tích hình chữ nhật là bao nhiêu mét vuông?
  \shortans{120}
  \loigiai{
    $15 \times 8 = 120\text{ m}^2$.
  }
\end{ex}

% Câu 5
\begin{ex}
  Một hình tam giác có đáy $14\text{ cm}$, chiều cao $10\text{ cm}$. Diện tích hình tam giác là bao nhiêu xăng-ti-mét vuông?
  \shortans{70}
  \loigiai{
    $(14 \times 10) : 2 = 70\text{ cm}^2$.
  }
\end{ex}

% Câu 6
\begin{ex}
  Một hình thang có đáy lớn $20\text{ cm}$, đáy bé $10\text{ cm}$, chiều cao $8\text{ cm}$. Diện tích hình thang là bao nhiêu xăng-ti-mét vuông?
  \shortans{120}
  \loigiai{
    $(20 + 10) \times 8 : 2 = 120\text{ cm}^2$.
  }
\end{ex}

% Câu 7
\begin{ex}
  Một hình tròn có bán kính $10\text{ cm}$. Diện tích hình tròn là bao nhiêu xăng-ti-mét vuông? (Lấy $\pi = 3{,}14$).
  \shortans{314}
  \loigiai{
    $3{,}14 \times 10 \times 10 = 314\text{ cm}^2$.
  }
\end{ex}

% Câu 8
\begin{ex}
  Một hình bình hành có đáy $16\text{ cm}$, chiều cao $9\text{ cm}$. Diện tích hình bình hành là bao nhiêu xăng-ti-mét vuông?
  \shortans{144}
  \loigiai{
    $16 \times 9 = 144\text{ cm}^2$.
  }
\end{ex}

% Câu 9
\begin{ex}
  Một hình thoi có hai đường chéo dài $14\text{ cm}$ và $10\text{ cm}$. Diện tích hình thoi là bao nhiêu xăng-ti-mét vuông?
  \shortans{70}
  \loigiai{
    $(14 \times 10) : 2 = 70\text{ cm}^2$.
  }
\end{ex}

% Câu 10
\begin{ex}
  Một mảnh vườn hình chữ nhật có diện tích $180\text{ m}^2$, chiều dài $15\text{ m}$. Chiều rộng mảnh vườn là bao nhiêu mét?
  \shortans{12}
  \loigiai{
    $180 : 15 = 12\text{ m}$.
  }
\end{ex}

% Câu 11
\begin{ex}
  Một tam giác có diện tích $54\text{ cm}^2$, chiều cao $9\text{ cm}$. Đáy tam giác dài bao nhiêu xăng-ti-mét?
  \shortans{12}
  \loigiai{
    $(54 \times 2) : 9 = 12\text{ cm}$.
  }
\end{ex}

% Câu 12
\begin{ex}
  Một hình thang có diện tích $96\text{ cm}^2$, chiều cao $8\text{ cm}$. Tổng hai đáy là bao nhiêu xăng-ti-mét?
  \shortans{24}
  \loigiai{
    $(96 \times 2) : 8 = 24\text{ cm}$.
  }
\end{ex}

% Câu 13
\begin{ex}
  Một thửa ruộng hình vuông có chu vi $120\text{ m}$. Diện tích thửa ruộng là bao nhiêu mét vuông?
  \shortans{900}
  \loigiai{
    Cạnh là $120 : 4 = 30\text{ m}$. Diện tích là $30 \times 30 = 900\text{ m}^2$.
  }
\end{ex}

% Câu 14
\begin{ex}
  Một hình chữ nhật có chu vi $80\text{ cm}$, chiều dài gấp 3 lần chiều rộng. Diện tích hình chữ nhật là bao nhiêu xăng-ti-mét vuông?
  \shortans{300}
  \loigiai{
    Nửa chu vi 40 cm. Rộng 10 cm, dài 30 cm. Diện tích: $30 \times 10 = 300\text{ cm}^2$.
  }
\end{ex}

% Câu 15
\begin{ex}
  Một bánh xe có đường kính $0{,}5\text{ m}$. Lăn 100 vòng được bao nhiêu mét? (Lấy $\pi = 3{,}14$).
  \shortans{157}
  \loigiai{
    Chu vi là $3{,}14 \times 0{,}5 = 1{,}57\text{ m}$. 100 vòng đi được $1{,}57 \times 100 = 157\text{ m}$.
  }
\end{ex}

% Câu 16
\begin{ex}
  Người ta rào một khu vườn hình chữ nhật dài $25\text{ m}$, rộng $15\text{ m}$ có cổng vào rộng $2\text{ m}$. Chiều dài hàng rào là bao nhiêu mét?
  \shortans{78}
  \loigiai{
    Chu vi là $(25 + 15) \times 2 = 80\text{ m}$. Hàng rào dài $80 - 2 = 78\text{ m}$.
  }
\end{ex}

% Câu 17
\begin{ex}
  Một căn phòng hình chữ nhật dài $8\text{ m}$, rộng $5\text{ m}$. Lát nền bằng các viên gạch vuông cạnh $50\text{ cm}$. Cần bao nhiêu viên gạch?
  \shortans{160}
  \loigiai{
    Diện tích phòng $40\text{ m}^2 = 400\,000\text{ cm}^2$. Gạch: $50 \times 50 = 2\,500\text{ cm}^2$. Số gạch: $400\,000 : 2\,500 = 160$ viên.
  }
\end{ex}

% Câu 18
\begin{ex}
  Một hình thoi có diện tích $48\text{ cm}^2$, một đường chéo là $12\text{ cm}$. Đường chéo kia là bao nhiêu xăng-ti-mét?
  \shortans{8}
  \loigiai{
    $(48 \times 2) : 12 = 8\text{ cm}$.
  }
\end{ex}

% Câu 19
\begin{ex}
  Một tam giác vuông có diện tích $24\text{ cm}^2$, một cạnh góc vuông $6\text{ cm}$. Cạnh góc vuông kia dài bao nhiêu xăng-ti-mét?
  \shortans{8}
  \loigiai{
    $(24 \times 2) : 6 = 8\text{ cm}$.
  }
\end{ex}

% Câu 20
\begin{ex}
  Nếu cạnh của hình vuông tăng lên gấp 3 lần thì chu vi hình vuông tăng lên gấp mấy lần?
  \shortans{3}
  \loigiai{
    Chu vi tỉ lệ thuận với cạnh nên tăng 3 lần.
  }
\end{ex}

% Câu 21
\begin{ex}
  Nếu cạnh của hình vuông tăng lên gấp 3 lần thì diện tích tăng lên gấp mấy lần?
  \shortans{9}
  \loigiai{
    Diện tích tăng $3 \times 3 = 9$ lần.
  }
\end{ex}

% Câu 22
\begin{ex}
  Một cái ao hình tròn bán kính $5\text{ m}$. Diện tích cái ao nhân đôi là bao nhiêu mét vuông? (Lấy $\pi = 3{,}14$).
  \shortans{157}
  \loigiai{
    Diện tích là $3{,}14 \times 25 = 78{,}5\text{ m}^2$. Nhân đôi được 157.
  }
\end{ex}

% Câu 23
\begin{ex}
  Hình thang có đáy lớn $18\text{ cm}$, đáy bé bằng một nửa đáy lớn, chiều cao $10\text{ cm}$. Diện tích hình thang là bao nhiêu xăng-ti-mét vuông?
  \shortans{135}
  \loigiai{
    Đáy bé là 9 cm. Diện tích là $(18 + 9) \times 10 : 2 = 135\text{ cm}^2$.
  }
\end{ex}

% Câu 24
\begin{ex}
  Một mảnh đất hình chữ nhật có chiều dài $30\text{ m}$, chiều rộng $20\text{ m}$. Người ta dành $20\%$ diện tích để làm lối đi. Diện tích lối đi là bao nhiêu mét vuông?
  \shortans{120}
  \loigiai{
    Diện tích là $30 \times 20 = 600\text{ m}^2$. Lối đi là $600 \times 0{,}2 = 120\text{ m}^2$.
  }
\end{ex}

% Câu 25
\begin{ex}
  Một khung tranh hình vuông có diện tích $81\text{ dm}^2$. Chu vi khung tranh là bao nhiêu đề-xi-mét?
  \shortans{36}
  \loigiai{
    Cạnh là $9\text{ dm}$. Chu vi là $9 \times 4 = 36\text{ dm}$.
  }
\end{ex}
